% !TeX root = ../深度学习与大语言模型技术原理.tex
% 由主文档通过 \input 载入；请编译主文档。
\setcounter{chapter}{1} % 保留原书章号。
\chapter{六个综合实验与关键自检}
\label{app:experiments}

本附录给出可执行的实验规格，属于学习任务设计，不表示这些实验已在本书编写过程中运行。模型规模和训练步数应根据可用算力选择，但基线、变量、划分和验收条件应在运行前明确。

\section{实验一：从零实现小型解码器}

输入为一份可合法使用的小型文本语料，输出为可训练和增量生成的模型。先固定分词器与训练、验证划分，实现嵌入、因果注意力、RoPE、RMSNorm、SwiGLU、残差和输出头，再按式~\eqref{eq:autoregressive} 训练。先验证很小数据上的损失能明显下降，再进行独立验证，避免将过拟合能力当成泛化能力。

验收包括：标签移位可手算，未来词元变化不影响过去输出，缓存与全前向 logits 在容差内一致，采样与结束规则明确。记录参数量、有效词元、困惑度、预填充与解码时间。将 MHA 改为 GQA 时，同时比较质量与缓存容量，不只展示生成案例。

\section{实验二：SFT、LoRA 与 QLoRA 对照}

选定一个规模可控的基座与领域指令数据，固定模板、回答掩码、数据划分和训练预算。建立未微调基线、LoRA 方案，并在硬件允许时加入全参数或 QLoRA 对照。检查多轮角色边界、空目标、截断、拼接掩码和累计词元归一化。

输出应包含领域质量、通用能力保持、可训练参数、峰值显存和训练速度。合并适配器后重新验证输出；量化合并若改变数值，应报告影响。不能以训练损失最低替代最终方案选择，也不能把未测的节省比例当成实际显存结果。

\section{实验三：奖励模型与直接偏好学习}

构造同提示下的偏好对，并用人工或独立标准抽检标签。先实现奖励差的逻辑损失，再实现式~\eqref{eq:dpo}。在极小数据上核对优选与拒绝回答互换后损失变化，确认参考模型冻结、提示不计分、序列对数概率正确相加。

对比 SFT 基线与 DPO 后模型的独立偏好、任务正确性、回答长度和目标似然。至少设置一个标签噪声或长度偏置对照，用于检验方法是在学习任务偏好还是利用表面特征。若加入 IPO 或 SimPO，应明确新的长度和间隔约定，避免把不同目标的损失数值直接比较。

\section{实验四：可验证任务上的组采样训练}

选择有独立验证条件的算术、代码或 SQL 任务。先构造评分器正常、边界和错误样例，测量误接受与误拒绝；SQL 可在隔离数据库中比较查询结果，还需考虑空结果、排序、重复行与只读约束。然后从固定策略对每题采样多回答，检查奖励分布和组内方差，再实现 RLOO 或 GRPO 更新。

报告任务成功率、全对组与全错组比例、平均长度、KL、裁剪比例、生成和训练耗时。设置只做拒绝筛选 SFT 的基线，以判断在线优化是否在相同预算下有额外价值。训练、奖励开发与最终评测题目应分开，不能通过修改评分器规则追逐测试成绩。

\section{实验五：可归因的检索增强生成}

准备带来源与相关性标注的小型文档库，依次构建 BM25、稠密检索、混合融合和重排。先相对精确向量搜索测 ANN 召回，再测语义相关性。随后构造统一证据格式与引用生成，逐项加入查询改写、上下文扩展或迭代检索。

报告 Recall@K、MRR 或 NDCG，以及回答正确性、引用支持和端到端延迟。加入人工理想证据实验，分离检索与生成错误。对无答案、矛盾来源、过期文档和权限过滤分别评测；只给一个漂亮问答演示不足以说明系统有效。

\section{实验六：受控软件工程智能体}

在隔离小仓库中提供检索、读取、补丁和测试工具，明确参数契约、可写范围、预算和停止条件。先以固定工作流完成任务，再加入动态规划；保存状态与工具结果，验证超时、重试、部分成功和循环检测。任何真实副作用都应通过任务已授权的受控接口发生。

独立评测包含任务完成率、补丁正确性、测试覆盖、无效动作数和成本。用同预算单智能体作为多智能体对照；若进一步收集轨迹训练，应把模型动作与环境观察分开，避免学习伪造工具输出。最终产物应是可审查补丁和验证证据，而非只含模型自述的成功日志。

\section{四个可手算的原理自检}

\begin{example}
\label{ex:unequal-microbatches}
两个微批量分别有2和8个有效词元，平均损失分别为1和3。比较“微批量平均的平均”与总体词元平均损失。
\end{example}
\begin{solve}
前者为 $(1+3)/2=2$，后者为 $(2\times1+8\times3)/10=2.6$。两者分别给微批量和词元等权；在变长回答训练中，误用前者会相对提高短回答的权重。对应依据是式~\eqref{eq:token-normalization}。
\end{solve}

\begin{example}
若草稿分布 $q=(0.8,0.2)$、目标分布 $p=(0.5,0.5)$，求单词元推测采样的接受率和拒绝后分布。
\end{example}
\begin{solve}
两个草稿词元的条件接受率分别为 $0.5/0.8=0.625$ 和1，总接受率为 $0.8\times0.625+0.2=0.7$。剩余质量为 $(0,0.3)$，所以拒绝后一定输出第二个词元。最终第一项概率为0.5，第二项为 $0.2+0.3=0.5$，与命题~\ref{prop:speculative-exact} 一致。
\end{solve}

\begin{example}
某题的四个回答奖励为 $(1,0,0,1)$。使用分母为组大小的方差定义，忽略稳定项，求 GRPO 优势；再求第一个回答的 RLOO 优势。
\end{example}
\begin{solve}
均值为 $1/2$、标准差为 $1/2$，因此 GRPO 优势为 $(1,-1,-1,1)$。第一个回答的留一基线为 $1/3$，RLOO 优势为 $2/3$。这说明两者的基线和缩放不同；不能仅因都使用组内奖励就认定数值目标相同。
\end{solve}

\begin{example}
某查询有两个相关文档，前四个结果的二值相关性为 $(0,1,0,1)$，求 Recall@4、Precision@4 与该查询的倒数排名。
\end{example}
\begin{solve}
两个相关文档均被召回，因此 Recall@4 为1，Precision@4 为 $2/4=1/2$。第一个相关结果位于第2名，倒数排名为 $1/2$。覆盖已完整，但头部排序仍有改善空间，这正是区分召回与重排的理由。
\end{solve}

